\section{Chemical Kinetics and Equilibrium}

\subsection{Reaction Rates}

\subsubsection{Stoichiometric reaction rate}
\[
a\ce{A}+b\ce{B}\rightarrow c\ce{C}+d\ce{D}
\]
\[
r=-\frac{1}{a}\frac{d[\ce{A}]}{dt}
=-\frac{1}{b}\frac{d[\ce{B}]}{dt}
=\frac{1}{c}\frac{d[\ce{C}]}{dt}
=\frac{1}{d}\frac{d[\ce{D}]}{dt}
\]
The reaction rate is normalized by stoichiometric coefficients so all species describe the same reaction progress.

\subsubsection{Average and instantaneous rate}
\[
r_{\mathrm{avg}}\approx -\frac{\Delta[\ce{A}]}{\Delta t}, \qquad
r=-\frac{d[\ce{A}]}{dt}
\]
The instantaneous rate is the slope of the concentration-time curve at a point.

\subsection{Rate Laws}

\subsubsection{General rate law}
\[
r=k[\ce{A}]^m[\ce{B}]^n
\]
The exponents are experimentally determined reaction orders. The overall order is \(m+n\).

\subsubsection{First-order integrated rate law}
\[
[\ce{A}]_t=[\ce{A}]_0e^{-kt}, \qquad
\ln[\ce{A}]_t=-kt+\ln[\ce{A}]_0
\]
A first-order plot of \(\ln[\ce{A}]\) versus \(t\) is linear with slope \(-k\).

\subsubsection{Second-order integrated rate law}
\[
\frac{1}{[\ce{A}]_t}=kt+\frac{1}{[\ce{A}]_0}
\]
A second-order plot of \(1/[\ce{A}]\) versus \(t\) is linear with slope \(k\).

\subsubsection{Half-life}
\[
t_{1/2}^{(1)}=\frac{\ln 2}{k}, \qquad
t_{1/2}^{(2)}=\frac{1}{k[\ce{A}]_0}
\]
First-order half-life is independent of initial concentration. Second-order half-life depends on \([\ce{A}]_0\).

\subsection{Temperature and Catalysis}

\subsubsection{Arrhenius equation}
\[
k=Ae^{-E_{\mathrm a}/RT}
\]
\(E_{\mathrm a}\) is activation energy and \(A\) is the frequency factor. Higher temperature increases the fraction of molecules above the activation barrier.

\subsubsection{Linear Arrhenius form}
\[
\ln k=-\frac{E_{\mathrm a}}{R}\frac{1}{T}+\ln A
\]
A plot of \(\ln k\) versus \(1/T\) has slope \(-E_{\mathrm a}/R\).

\subsubsection{Two-temperature Arrhenius relation}
\[
\ln\left(\frac{k_2}{k_1}\right)
=-\frac{E_{\mathrm a}}{R}\left(\frac{1}{T_2}-\frac{1}{T_1}\right)
\]
This relation compares rate constants at two temperatures when \(E_{\mathrm a}\) is constant.

\subsubsection{Catalyst effect}
\[
E_{\mathrm a,cat}<E_{\mathrm a,uncat}
\]
A catalyst lowers the activation barrier and increases rate. It does not change the thermodynamic equilibrium constant.

\subsection{Equilibrium}

\subsubsection{Reaction quotient}
\[
\alpha\ce{A}+\beta\ce{B}\rightleftharpoons\gamma\ce{C}+\delta\ce{D}
\]
\[
Q=\frac{a_{\ce{C}}^\gamma a_{\ce{D}}^\delta}{a_{\ce{A}}^\alpha a_{\ce{B}}^\beta}
\]
Activities are usually represented by concentrations for solutes and partial pressures for gases in this course. Pure solids, pure liquids, and solvents are omitted.

\subsubsection{Equilibrium constant}
\[
K=Q_{\mathrm{eq}}
\]
At equilibrium, forward and reverse rates balance and the reaction quotient equals \(K\).

\subsubsection{Gibbs free energy and reaction quotient}
\[
\Delta G=\Delta G^\circ+RT\ln Q
\]
Non-standard spontaneity is determined by \(\Delta G\), not by \(\Delta G^\circ\) alone.

\subsubsection{Standard free energy and equilibrium constant}
\[
\Delta G^\circ=-RT\ln K, \qquad K=e^{-\Delta G^\circ/RT}
\]
\(K>1\) corresponds to \(\Delta G^\circ<0\) and product-favored equilibrium.

\subsubsection{Le Chatelier's principle}
\[
\text{disturbance} \Rightarrow \text{shift opposing the disturbance}
\]
Changes in concentration, pressure, volume, or temperature shift an equilibrium to reduce the imposed change.

\subsubsection{Gas equilibrium constant}
\[
K_p=\frac{(p_{\ce{C}})^\gamma(p_{\ce{D}})^\delta}{(p_{\ce{A}})^\alpha(p_{\ce{B}})^\beta}
\]
For gas equilibria, partial pressures are used in the same stoichiometric powers as concentration expressions.
